% SGA 3, Expose VII_B, English reconstruction -- Loop 1.
% Scope: Section 5.1.2 from its opening through the filtration lemma
% statement on authority local page 84; the proof follows in component 78.
% Combined-reader page: 597. Re-edition printed page: 586.
% Original-edition marginal locators in scope: 553--554.
% Editorial notes in scope: 146--147.
% Authority: Polo--Gille Exp7B-13oct24.pdf,
% SHA-256 D303690E0195654BBE7C3DAA5D89EA23E24CD709A3575217234AC19AE8836694.
% Jacob Reinhold's CC BY 4.0 Markdown was consulted as comparison lineage.
% No genuine diagram occurs in this component; all mathematics is native LaTeX.

\subsubsection*{5.1.2}
\phantomsection
\label{sga3:VIIB:5.1.2}

用 \(\mathbf{Alpg}/k\) 表示\emph{分次预有限 \(k\)-代数}的范畴。此范畴的对象由一列伪紧 \(k\)-模
\[
  A_0,A_1,\ldots,A_n,\ldots
\]
及
\[
  \prod_{n\geq0}A_n
\]
上的预有限代数结构组成，满足
\[
  A_nA_m\subset A_{m+n}.
\]
这里借助典范单射将 \(A_n\) 与 \(\prod_{i\geq0}A_i\) 的一个子集等同。态射
\[
  \psi:(A_n)\longrightarrow(B_n)
\]
是一列连续线性映射
\[
  \psi_n:A_n\longrightarrow B_n
\]
使得对 \(a\in A_m\) 和 \(a'\in A_n\)，有
\[
  \psi_{m+n}(aa')=\psi_m(a)\psi_n(a')
\]

\medskip
\noindent\textit{定义。}
两个分次预有限 \(k\)-代数 \((A_n)\) 和 \((B_n)\) 显然在 \(\mathbf{Alpg}/k\) 中有余积\footnote{%
\label{sga3:VIIB:nde:146}编者将“直和”换成“余积”。}。其第 \(n\) 个分量为
\[
  \prod_{i=0}^{n}
  A_i\widehat\otimes_k B_{n-i},
\]
即伪紧 \(k\)-模
\[
  A_i\widehat\otimes_k B_{n-i}.
\]
的积。此余积记作
\[
  (A_n)\widehat\otimes_k(B_n).
\]

因此，\(\mathbf{Alpg}/k\) 的对象 \((A_n)\) 上的余群结构，是指一组连续 \(k\)-线性映射
\[
  \Delta_n:A_n\longrightarrow
  \prod_{i=0}^{n}A_i\widehat\otimes_k A_{n-i},
  \qquad
  \varepsilon:A_0\longrightarrow k,
\]
它们在对 \(i\geq1\) 置 \(\varepsilon(A_i)=0\) 后，在 \(\prod_{n\geq0}A_n\) 上诱导 \(\mathbf{Alp}/k\) 中的余群结构。

最后，对 \(\mathcal C\) 的每个对象 \((A,J)\)，用
\[
  \operatorname{Gr}_J(A)
\]
表示与 \(A\) 按 \(J\) 的幂的闭包 \(\overline{J^n}\) 所得的滤过相联系的分次预有限代数。因此，
\[
  \operatorname{Gr}_J(A)_n
  =
  \overline{J^n}/\overline{J^{n+1}},
\]
且 \(\operatorname{Gr}_J(A)\) 的乘法由 \(A\) 的乘法诱导。

\medskip
\noindent\textbf{引理。}
\phantomsection
\label{sga3:VIIB:5.1.2:filtration-lemma}
\footnote{\label{sga3:VIIB:nde:147}原文在 \(k\) 为域时陈述此引理，并将此时的证明留给读者。}
{\itshape 设 \(U,V\) 为两个伪紧 \(k\)-模，\(U\) 拓扑平坦，并设
\[
  U=U_0\supset U_1\supset\cdots,
  \qquad
  V=V_0\supset V_1\supset\cdots
\]
为两个闭 \(k\)-子模的递减序列。用闭子模
\[
  W_n
  =
  U_n\widehat\otimes_kV_0
  +
  U_{n-1}\widehat\otimes_kV_1
  +\cdots+
  U_0\widehat\otimes_kV_n.
\]
赋予完备化张量积
\[
  W=U\widehat\otimes_kV
\]
滤过。假设每个 \(U_i/U_{i+1}\) 都在 \(k\) 上拓扑平坦，从而对每个 \(n\)，\(U/U_n\) 以及 \(U_n\) 也拓扑平坦。则对每个 \(n\)，有同构
\[
  W_n/W_{n+1}
  \simeq
  \bigoplus_{i+j=n}
  (U_i/U_{i+1})\widehat\otimes_k(V_j/V_{j+1}).
\]}
